% !TeX root = ../../金融经济学_熊和平_习题.tex
% 源：重拍 113800（印张 19）、113807（印张 20）；225493.pdf p.1

\zhchapter{基本框架及分析模式}

\begin{tcolorbox}[summary,title=本章小结]
金融经济学在分析问题时需要建立一定的模型，这种模型即是描述问题所涉及的基本框架。经济框架一般涉及三个方面：其一是经济环境；其二是证券市场；其三是经济人。

本书采用简单的单期框架，主要关注两个时点：期初投资和期末获得收益。经济环境还包括信息结构，即未来世界所处的状态及其分布。本书强调证券市场这一核心概念。经济人包括个人、厂商和政府；家庭金融把家庭从个人中分离出来单独研究。

一般均衡分析分两步：第一步，在偏好、禀赋和市场结构的约束下，经济人进行优化，确定消费或资产需求与价格之间的函数关系；第二步，利用市场出清条件（总禀赋等于总需求）联立求解均衡价格与配置。无套利分析法通过比较证券与组合的支付向量，寻找复制关系，从而对目标证券定价；其运用要求证券可被复制，并通常假设至少一位经济人具有不满足性、套利机会可被迅速利用（本书第六章展开具体方法）。

纯交换经济下证券供给外生给定，配置的可行性与 Pareto 最优（定义 1.2，约束 (1.15)--(1.16)）是评价资源配置的重要概念；引入生产后即为生产性经济，第十三章专门讨论。
\end{tcolorbox}

\xhpconceptitems{
    \item 经济环境，信息结构，事件，样本空间
    \item 投资组合，支付向量，支付空间，市场结构
    \item 一般均衡分析，套利分析，Lucas 树
    \item 风险证券，无风险证券，市场化，禀赋，消费计划，经济人
    \item 相对定价法，绝对定价法，市场出清，配置，Pareto 最优配置
}

% !TeX root = ../../金融经济学_熊和平_习题.tex

\begin{problem}
\problemno{1.1} 金融经济学是从哪几个方面来描述经济环境的？
\end{problem}
\begin{problem}
\problemno{1.2} 什么是纯交换经济？什么是生产性经济？两者的区别和联系是什么？
\end{problem}
\begin{problem}
\problemno{1.3} 金融经济学从哪些方面描述经济人？
\end{problem}
\begin{problem}
\problemno{1.4} 什么叫市场结构？
\end{problem}
\begin{problem}
\problemno{1.5} 什么是投资组合？投资组合一般有哪几种描述方式。
\end{problem}
\begin{problem}
\problemno{1.6} 什么是一般均衡分析方法？一般均衡分析方法一般包括哪几步？
\end{problem}
\begin{problem}
\problemno{1.7} 什么是无套利分析法？请简述该方法的基本原理。
\end{problem}

